\documentclass[11pt,letterpaper]{article}
\usepackage[T1]{fontenc}
\usepackage{lmodern,amsmath,amssymb,amsthm,mathtools}
\usepackage[margin=0.87in]{geometry}
\usepackage{booktabs,longtable,array,microtype,xcolor,enumitem,fancyhdr}
\usepackage{hyperref,xurl}
\hypersetup{colorlinks=true,linkcolor=blue!40!black,citecolor=blue!40!black,
urlcolor=blue!40!black,pdftitle={Pooling exact power certificates for Diophantine consistency},
pdfauthor={AI-assisted research prepared for Vladimir Reshetnikov}}
\newtheorem{theorem}{Theorem}[section]
\newtheorem{lemma}[theorem]{Lemma}
\newtheorem{proposition}[theorem]{Proposition}
\newtheorem{corollary}[theorem]{Corollary}
\theoremstyle{definition}\newtheorem{definition}[theorem]{Definition}
\theoremstyle{remark}\newtheorem{remark}[theorem]{Remark}
\newcommand{\Z}{\mathbb Z}\newcommand{\N}{\mathbb N}
\newcommand{\Con}{\operatorname{Con}}
\newcommand{\G}{\mathcal G}
\newcommand{\Pow}{\mathcal P}
\setlength{\emergencystretch}{2em}
\setlist{itemsep=3pt,topsep=5pt}
\pagestyle{fancy}\fancyhf{}
\fancyhead[L]{\small Diophantine consistency}
\fancyhead[R]{\small Pooled exact powers}
\fancyfoot[C]{\thepage}\setlength{\headheight}{14pt}
% Generated by code/export_equation.py.
\newcommand{\NumVars}{225}
\newcommand{\LastVar}{224}
\newcommand{\NumRows}{162}
\newcommand{\LastRow}{161}
\newcommand{\NumTerms}{7,867}
\newcommand{\CoeffHeight}{128,446}

\title{Pooling exact power certificates for\\Diophantine consistency\\[6pt]
\large A fixed quartic in \NumVars\ unrestricted integer variables}
\author{AI-assisted research prepared for Vladimir Reshetnikov\\
\small ProveIt research program}
\date{25 October 2026}
\begin{document}
\maketitle
\begin{abstract}
We construct a fixed polynomial of total degree four in \NumVars\
unrestricted integer variables whose integer unsolvability is provably
equivalent in ZFC to \(\Con(\mathrm{ZFC})\), by composing the new arithmetic
proof with the repository's identified C222 refutation calibration.
The polynomial is the sum of squares of \NumRows\ literal quadratic
residuals; its collected expansion has \NumTerms\ nonzero monomials and
coefficient height \CoeffHeight. The inspected parent has 245 variables,
181 residuals, 5,304 monomials, and the same degree and height.
The principal new construction replaces all power congruences sharing one
base by two quadratic equations with two auxiliary integers, independently
of the number of requests. A sum of squared errors, a signed Pell
congruence, and the already available square frame prove simultaneous
exactness, including exclusion of negative exponents at nonunit bases.
The modulus itself need not be a coordinate. Eleven requests at one base
and three at another save twelve variables. Cofinal strict ranges and a
smaller beta base save five more; exact changes of auxiliary coordinates
save three. The delivered equation encodes a numerical proof calculus
directly and uses neither K215 nor Jones's universal polynomial.
The result improves this explicit direct quartic family. It is not a
global minimality theorem or a new end-to-end proof-assistant certificate.
\end{abstract}
\clearpage
\tableofcontents
\clearpage

\section{The equation, theorem, and comparison}

Let \(\mathbf x=(x_0,\ldots,x_{\LastVar})\in\Z^{\NumVars}\). The equation is
\begin{equation}\label{eq:main}
\boxed{\qquad P_{\NumVars}(\mathbf x)
       =\sum_{j=0}^{\LastRow}R_j(\mathbf x)^2=0.\qquad}
\end{equation}
Appendix~\ref{app:residuals} prints every residual \(R_j\), with every
coefficient and coordinate index. The same literal specification is
\path{equation/equation.json}; the collected polynomial is
\path{equation/expanded_polynomial.json}. All coefficients are fixed
integers. No factorial, binomial operator, variable exponent, unspecified
parameter, or black-box universal polynomial remains in
\eqref{eq:main}.

\begin{theorem}\label{thm:main}
For the literal polynomial in \eqref{eq:main}, the arguments below, composed
with the retained arithmetic and C222/ZF/ZFC refutation calibrations, prove
\[
\mathrm{ZFC}\ \vdash\
\left(\forall\mathbf x\in\Z^{\NumVars},\ P_{\NumVars}(\mathbf x)\ne0\right)
\quad\longleftrightarrow\quad\Con(\mathrm{ZFC}).
\]
The polynomial has exact total degree four, \NumTerms\ collected nonzero
monomials, and maximum absolute coefficient \CoeffHeight.
\end{theorem}

The consistency assertion is syntactic consistency for the effective ZFC
presentation fixed by the inherited calibration. The construction assumes
neither consistency nor inconsistency. Its coefficients and verification
artifacts are obtained by finite arithmetic independently of whether a
refutation exists.

The comparison is pinned to ProveIt commit
\begin{center}\small\texttt{a0ef2df5bb1f73c35cd334bbc7f53553350ef06a}.\end{center}
All nine incoming archives at that snapshot were downloaded and their
relevant reports inspected. The immediate baseline is the
\emph{Joint Norm} archive, whose 245-variable direct polynomial supersedes
the older 252-variable endpoint for this comparison \cite{repo,joint}.
The main collection's top-level summary does not yet include that incoming
improvement. It would therefore be misleading to compare only with its
older 282-variable direct result.

\begin{center}
\begin{tabular}{lrrr}
\toprule
Measure & Inspected parent & This equation & Change\\
\midrule
Unrestricted integer variables & 245 & \NumVars & $-20$\\
Quadratic residuals & 181 & \NumRows & $-19$\\
Total degree & 4 & 4 & 0\\
Collected monomials & 5,304 & \NumTerms & $+2,563$\\
Coefficient height & 128,446 & \CoeffHeight & 0\\
Largest residual coefficient & 256 & 256 & 0\\
\bottomrule
\end{tabular}
\end{center}

The reduction in variables is approximately \(8.16\%\). The larger
monomial count is a real tradeoff: squaring a pooled residual creates more
cross terms. The claim is fewer variables at the same degree and
coefficient height, not improvement of every complexity measure.

\paragraph{What ``simplest'' can mean.}
The repository also contains equations with far fewer variables but
substantially larger degree or large fixed coefficient encodings
\cite{forms}. Sun's ten-integer undecidability theorem concerns a different
degree/variable tradeoff \cite{sun}; the present construction does not
improve that general variable bound. Our objective here is a direct
quartic with moderate literal coefficients and a transparent arithmetic
certificate. We do not know an absolute simplest equation, and no
literature-wide priority or minimality claim is made.

\section{The new idea: two equations for many exact powers}
\label{sec:pool}

The inherited construction already classifies a finite family of signed
Pell pairs at one large parameter. Its remaining power checks had one
congruence quotient per request and one modulus per base. Those separate
quotients can be removed in groups. The important feature is that the
aggregation remains quadratic. The earlier batched-certificate proposal
shares the Pell environment while keeping individual power congruences
\cite{batched}; the present pooling removes those individual quotients.

\subsection{Signed Pell notation and the classical congruence}

For an integer \(A\) with \(|A|>1\), define the Pell polynomials by
\[
\chi_A(0)=1,\qquad\psi_A(0)=0,
\]
\[
\begin{pmatrix}\chi_A(k+1)\\\psi_A(k+1)\end{pmatrix}
=\begin{pmatrix}A&A^2-1\\1&A\end{pmatrix}
 \begin{pmatrix}\chi_A(k)\\\psi_A(k)\end{pmatrix}.
\]
The matrix has determinant one, so the definition extends to every
\(k\in\Z\), with \(\chi_A(-k)=\chi_A(k)\) and
\(\psi_A(-k)=-\psi_A(k)\). In particular,
\[
\chi_A(k)^2-(A^2-1)\psi_A(k)^2=1.
\]
These functions occur in the proof only. Existing quadratic equations
certify their coordinates in the delivered polynomial.

For an integer base \(n\), put
\begin{equation}\label{eq:modulus}
T=2An-n^2-1.
\end{equation}
Then \(\gcd(n,T)=1\), since \(T\equiv-1\pmod n\). The standard power
congruence is
\begin{equation}\label{eq:pellpower}
\chi_A(k)+(n-A)\psi_A(k)\equiv n^k\pmod T,
\qquad k\in\Z,
\end{equation}
where a negative exponent on the right is interpreted in the group of
units of \(\Z/(T)\). To prove it, substitute \(n-A\) for the formal
square root of \(A^2-1\):
\((n-A)^2\equiv A^2-1\pmod T\), and the fundamental unit maps to
\(A+(n-A)=n\). Positive powers follow by multiplication and negative
powers by invertibility. This proves the signed version directly.
Classical Diophantine power and factorial representations are also
formalized in existing libraries and work such as \cite{pak}; our new
step is the following aggregate use of the congruence.

\subsection{An exact quadratic identity}

Suppose a group has one common base \(n\) and requests
\((k_i,m_i)\), \(1\le i\le r\). Write
\[
X_i=\chi_A(k_i),\qquad Y_i=\psi_A(k_i),\qquad Z_i=AY_i.
\]
The \(Z_i\) are existing coordinates, with quadratic definitions.
For the proof abbreviate
\[
D_i=m_i-X_i+Z_i,\qquad
H=\sum_{i=1}^rY_i(m_i-X_i),\qquad
S=\sum_{i=1}^r(D_i^2-Y_i^2).
\]
Introduce only two integers \(K,L\), and impose
\begin{align}
K-2\sum_{i=1}^rY_i(m_i-X_i)-(2A-n)L&=0,
                                                        \label{eq:pool1}\\
\sum_{i=1}^r\bigl((m_i-X_i+Z_i)^2-Y_i^2\bigr)-nK+L&=0.
                                                        \label{eq:pool2}
\end{align}
Both residuals have total degree at most two. Neither \(H\), \(S\),
\(D_i\), nor \(T\) is an additional coordinate. The number of auxiliary
variables is independent of \(r\).

\begin{lemma}[Exact aggregate identity]\label{lem:identity}
Assume \(Z_i=AY_i\). Equations~\eqref{eq:pool1}--\eqref{eq:pool2}
imply
\begin{equation}\label{eq:erroridentity}
\sum_{i=1}^r\bigl(m_i-X_i+(A-n)Y_i\bigr)^2
=T\left(L-\sum_{i=1}^rY_i^2\right).
\end{equation}
Conversely, if each individual power error on the left is divisible by
\(T\ne0\), integers \(K,L\) satisfying the two quadratic rows exist.
\end{lemma}

\begin{proof}
Substitution of \(K=2H+(2A-n)L\) into the second row gives
\(S-2nH=TL\). On the other hand,
\begin{align*}
\sum_i(D_i-nY_i)^2
 &=\sum_iD_i^2-2nH-(2An-n^2)\sum_iY_i^2\\
 &=S-2nH-T\sum_iY_i^2.
\end{align*}
This is \eqref{eq:erroridentity}. Conversely set
\[
L=\sum_iY_i^2+
 \frac{\sum_i(m_i-X_i+(A-n)Y_i)^2}{T},
\qquad K=2H+(2A-n)L.
\]
The assumed divisibilities make \(L\) integral, and the identity proves
both rows. No sign condition on either new integer is needed.
\end{proof}

One cannot infer that each error vanishes merely from
\(T\mid\sum e_i^2\). The next theorem supplies the essential small-error
argument. It applies to the small semantic errors obtained after Pell
reduction, not to the very large raw Pell coordinates.

\subsection{The existing square frame is sufficient}

The inherited frame is
\begin{equation}\label{eq:frame}
W=2+\sum_{a\in\mathcal A}a^2,
\qquad |A|>W^W,
\end{equation}
where every actual base, requested exponent, and proposed output is an
atom in \(\mathcal A\). Repeated atoms cause no problem. In our application
\(W\ge258\). Thus the stronger bounds
\(n^2,k_i^2,m_i^2<W\), not just \(|n|,|k_i|,|m_i|<W\), are available.

\begin{theorem}[Quadratic pooling of signed powers]\label{thm:pool}
Suppose \(W\ge258\), \(1\le r\le15\), \(n\ne0\),
\[
|A|>W^W,\qquad n^2,k_i^2,m_i^2<W,
\]
and the pairs \((X_i,Y_i)\) are the canonical signed Pell pairs at the
indices \(k_i\), with \(Z_i=AY_i\). Then the two quadratic equations
\eqref{eq:pool1}--\eqref{eq:pool2} force
\[
\begin{cases}
k_i\ge0\ \text{and}\ m_i=n^{k_i},& |n|\ge2,\\
m_i=n^{k_i},& n\in\{-1,1\}.
\end{cases}
\]
For unit bases the integer power is defined for every signed index.
Conversely, these exact power values admit \(K,L\).
\end{theorem}

\begin{proof}
Put \(b=\lfloor\sqrt{W-1}\rfloor\). Then \(b\ge16\),
\(|n|,|k_i|,|m_i|\le b\), and \(b+k_i\ge0\). By
\eqref{eq:erroridentity} and \eqref{eq:pellpower}, multiplication by
\(n^{2b}\) gives
\[
T\mid E,\qquad
E:=\sum_{i=1}^r\bigl(m_i n^b-n^{b+k_i}\bigr)^2.
\]
This is an ordinary nonnegative integer: the exponent shift has removed
every inverse. As \(b\ge2\),
\[
|m_i n^b-n^{b+k_i}|
 \le b^{b+1}+b^{2b}\le2b^{2b}.
\]
Consequently,
\begin{equation}\label{eq:smallerror}
0\le E\le4r b^{4b}<W^{2b+1}\le W^W.
\end{equation}
Here \(4r\le60<W\), \(b^2<W\), and
\(2b+1\le b^2+1\le W\). Also
\[
|T|\ge2|A||n|-n^2-1\ge2|A|-W>|A|>W^W.
\]
Thus \(0\le E<|T|\) and \(T\mid E\), so \(E=0\). Every summand
vanishes. If \(k_i\ge0\), cancel the nonzero integer \(n^b\) to obtain
\(m_i=n^{k_i}\). If \(k_i=-j<0\), cancellation gives
\(m_i n^j=1\). This is impossible for \(|n|\ge2\); for \(n=\pm1\)
it gives exactly the asserted unit power. Conversely the classical
congruence makes every individual error divisible by \(T\), and
Lemma~\ref{lem:identity} constructs \(K,L\).
\end{proof}

The number fifteen is a convenient bound for this particular application.
The identical proof works for any number \(r\) of requests when
\(4r<W\) and the other frame hypotheses hold. The two auxiliary
coordinates do not grow with the number of requests.

The use of the square bounds is decisive. A coarser estimate using only
\(|k_i|<W\) would suggest enlarging the Pell frame. That enlargement is
unnecessary here. The final polynomial retains the original odd Pell-frame
certificate unchanged. The finite checker reconstructs the frame as a sum
of 36 affine squares after the removed request and authenticates every
pooled base, exponent, and output against that list.

\subsection{The two actual groups and the variable saving}

After the strict-horizon change below, the direct source has seventeen
power requests. Fourteen belong to the two groups in this table.
The symbols are semantic values, before the final integral coordinate
changes.

\begin{center}
\begin{tabular}{lll}
\toprule
Base & Requested indices & Number\\
\midrule
\(q+1\) & \(F^2+t,a,b,c,d,S,S',u,v,C,h\) & 11\\
\(q\) & \(t,y+1,x\) & 3\\
\(F^2+1\) & \(t\) & 1\\
\(B\) & \(x\) & 1\\
\(BH\) & \(x\) & 1\\
\bottomrule
\end{tabular}
\end{center}

The nine independent CRT representatives are
\(a,b,c,d,S,S',u,v,C\). The derived argument \(h\) is a fixed
constructor polynomial in the four schema representatives. The last
three requests remain individual, which makes the bootstrap simple.
Each pooled group removes its old modulus coordinate and all its old
quotients, and introduces two integers. The savings are therefore
\[
(11+1)-2=10,\qquad(3+1)-2=2.
\]
The existing Pell pairs and their multiplication coordinates are retained.
The theorem is not using unrecorded powers as primitives; the entire
certificate consists of the retained quadratic rows and the four new
quadratic rows.

\section{Cofinal horizons remove a complete power request}
\label{sec:strict}

The horizon is selected existentially from
\(t=16+z^2\). A finite refutation permits arbitrarily large choices of
this form before the factorial and accumulator states are constructed.
This freedom makes weak upper bounds unnecessarily expensive.

\subsection{Strict local ranges}

The parent uses binomial divisibilities with lower index \(t+1\) for
four local schema parameters and one locally decoded derived-label sum.
Use lower index \(t\) instead. They now enforce
\[
0\le a,b,c,d<t,\qquad 0\le J(a,b,c,d)<t
\]
on every decoded local tuple. All inherited soundness arguments require
only the weaker bounds \(\le t\), so the stronger conditions suffice.
The one-sided prime-power binomial range lemma gives this conclusion
in exactly the same manner as in the parent.

The notation must be kept straight. The local value \(J(a,b,c,d)\) is
bounded by \(t\). The value \(h=J(A_a,A_b,A_c,A_d)\) on the large
global CRT representatives is generally much larger and is the semantic
Pell anchor. No bound \(h<t\) is imposed.

For completeness, take a finite proof tree, list all its parameter tuples,
and add the fixed truth-padding tuple. Choose \(t=16+z^2\) strictly larger
than its number of expansion steps, every parameter, and every full
derived-label sum on that finite list. These sums are independent of
\(F=t!\), so this choice is not circular.

At the parent's precision \(2x\), a global representative is congruent
to the intended local value \(w<t\) modulo \(p^{2xm}\). If the global
value is smaller than \(t\), the required binomial is zero. Otherwise
its falling factorial has a factor congruent to zero modulo
\(p^{2xm}\); division by \(t!\) loses at most \(t-1<xm\) powers of
\(p\). Thus the required \(p^{xm}\) divisibility holds. Polynomial
congruence proves the same statement for the derived-label value.

\subsection{One denominator disappears}

Put \(M=q^x\). The parent normalizes digit remainders by ratios
\(M/q^t\), \(M/q^{t+1}\), and \(M/q^{y+1}\), with a fourth group
already using \(M\). After tightening the five ranges, every use of
\(M/q^{t+1}\) can use \(M/q^t\). Delete the output \(q^{t+1}\), its
power quotient, the private wire for its one-step Pell addition, and its
normalization ratio. This saves four variables and three rows.

The removed output's square is also deleted from the square frame.
Every retained base, exponent, and output remains bounded by an explicit
square. The unchanged \(t\) pair is still needed by both \(q^t\) and
\((F^2+1)^t\). The modified early bootstrap obtains \(x\ge t\), rather
than \(x\ge t+1\); Section~\ref{sec:bootstrap} proves the adjusted
estimates explicitly.

\section{A smaller beta base}
\label{sec:beta}

Let \(V=(F^2+1)^t\). Replace the parent's
\(y=V(V-1)\) by the affine identity
\begin{equation}\label{eq:beta}
y=V-1.
\end{equation}
This is another use of the cofinal horizon, and eliminates one coordinate.

Starting with one pending formula, each backward expansion of a proof
tree replaces one formula by at most two premises. If the finite tree
requires \(L\) expansions, every pending multiset has at most
\(L+1\le t\) elements. Truth padding adds only a weight-one operation
and does not enlarge the nontrivial state. Every formula label \(j\)
lies below \(t\), and the inherited multiplicative weight is
\(1+jF\). Since \(F=t!\ge t\),
\[
1+jF\le1+(t-1)F<F^2+1.
\]
Every positive accumulator state is therefore strictly below
\((F^2+1)^t=V\), and at most \(y\). Also \(F\mid V-1=y\).
This is exactly the divisibility needed to make the beta moduli
\(1+(i+1)y\), for \(i=0,\ldots,t\), pairwise coprime. Indeed a common
prime divisor would divide an index difference, hence \(F\), hence
\(y\), which is incompatible with dividing \(1+(i+1)y\).

After factorial recovery, \(y>F+2\) still holds by an enormous margin.
Thus the endpoint and amplification arguments retain their hypotheses.
Before that recovery, \(V\ge1\) gives \(y\ge0\), including the
temporary branch \(F=0,V=1,y=0\), which is separately excluded below.

There is an offset in the centered parent coordinates: actual
\(y=x_2-1\), whereas \(V=x_6\). The literal new row is consequently
\(x_2-x_6=0\), not \(x_2-x_6+1=0\). Its exact elimination is checked
in the supplied certificate. This detail matters for both endpoints.

\section{Soundness in an acyclic order}
\label{sec:bootstrap}

We give the dependency order in some detail because signed witnesses and
the deleted initial guards make premature positivity assumptions unsafe.
The argument retains the parent's rounded factorial certificate and
shared remainder norm; the changed pieces are stated explicitly.

\subsection{What is available first}

The retained signed Pell cap and pooled coordinate-orientation argument
give all fourteen principal pairs at their requested signed indices,
before interpreting any front power. They use the square frame, a
separate squared semantic anchor, and bounded residues modulo \(A-1\).
Their hypotheses are unchanged by pooling; the cap does not depend on
individual power quotients. The retained proof texts are listed in
Section~\ref{sec:calibration}.

Among the three individual power requests, the request at \(F^2+1\)
first gives \(V=(F^2+1)^t\), since its index \(t\ge16\) is already known.
The shared norm has the semantic form
\begin{equation}\label{eq:jointnorm}
M^2=1+4\left(e^2+\sum_{j=1}^{11}\rho_j^2\right)+\G(\alpha,\beta,\gamma),
\qquad \G(u,v,w)=u^2+v^2+w(w+1),
\end{equation}
and the ceiling rows are
\begin{equation}\label{eq:ceiling}
Bd=M,\qquad FM-dV-e=0.
\end{equation}
Here \(\G\ge0\) for all integer inputs. Thus \(M,B,d\ne0\), and
\(e=d(FB-V)\) gives
\begin{equation}\label{eq:halfwidth}
|FB-V|<|B|/2.
\end{equation}
Write \(b_0=|B|\), \(f=|F|\). It follows that
\(1\le V<(f+1/2)b_0\). If \(f=0\), then \(b_0>2\). If
\(f\ge1\) and \(b_0\le f\), then
\((f^2+1)^t\ge(f^2+1)^2>f^2+f/2\), a contradiction. Hence
\begin{equation}\label{eq:frontmagnitudes}
b_0\ge2,\qquad |F|<b_0,\qquad V<b_0^2.
\end{equation}

\subsection{The unit-base case licenses the first pooled group}

Set \(H=1+2\sum a_i\), where the sum runs over the nine actual CRT
representatives. It is odd, hence nonzero, without any sign premise.
The individual requests at \(B\) and \(BH\) force
\[
x\ge0,\qquad N=B^x,\qquad q=(BH)^x.
\]
In particular \(q\ne0\), so Theorem~\ref{thm:pool} applies to the
three requests at base \(q\), even if temporarily \(q=1\).
If \(x=0\), then \(q=1\), and the pooled request for \(M=q^x\)
gives \(M=1\). This contradicts \(Bd=M\) and \(|B|\ge2\).
Therefore \(x\ge1\) and \(|q|\ge2\).

The retained denominator identity \(q^t r_t=q^x\) now gives
\(x\ge t\): a negative power of an integer of absolute value at least
two cannot be an integer ratio. Consequently
\[
|q|\ge b_0^x\ge2^{16},\qquad |q+1|\ge2^{16}-1\ge2.
\]
The eleven-request group at base \(q+1\) now forces every actual
representative to be nonnegative and all its powers to be exact. It also
forces the derived requested index to be nonnegative; the independent
anchor has already excluded zero. Thus \(H\ge1\).

This order explains why the theorem's unit-base clause matters. Demanding
\(|q|\ge2\) as a premise before proving \(x\ne0\) would leave a gap.

\subsection{Factorial digit recovery, including the zero candidate}

Let \(Q=|q|\). The binomial numerator is now the exact signed-radix
power \((q+1)^{F^2+t}\). If \(F\ne0\), each coefficient through degree
\(t\) is at most \(V=(F^2+1)^t\). The signed low tail \(r^0\) at
position \(t\) therefore satisfies
\[
\frac{|r^0|}{Q^t}<\frac{V}{Q-1}
 <\frac{2}{b_0^{x-2}}\le\frac{2}{2^{14}}<\frac1{32}.
\]
The common norm bounds the absolute value of every chosen remainder by
half its absolute modulus. Two congruent remainders, one of absolute
value less than half that modulus and the canonical one of absolute
value less than one thirty-second of it, cannot differ by a nonzero
multiple of the modulus.
They are equal. The next digit gives
\[
B\equiv\binom{F^2+t}{t}\pmod q.
\]
Both values have difference less than
\(b_0+b_0^2<Q\), so \(B=\binom{F^2+t}{t}>0\).

When \(F=0\), it would be wrong to bound all \(\binom tj\) by one.
Instead \(V=1\), \eqref{eq:halfwidth} gives \(b_0\ge3\), and
\[
\frac{|r^0|}{Q^t}<\frac{2^t}{Q-1}
 <2(2/3)^t\le2(2/3)^{16}<\frac1{32}.
\]
The same recovery forces \(B=\binom tt=1\), contradicting \(b_0\ge3\).
Thus this branch has no zero.

\subsection{The rounded factorial lemma}

For completeness we recall the inherited elementary fixed-point argument
on the range needed here. If \(t\ge16\), \(F\in\Z\),
\[
B=\binom{F^2+t}{t},\qquad V=(F^2+1)^t,
\]
then \(|F-V/B|<1/2\) holds exactly when \(F=t!\). Put
\(L_0=t!\), \(a=t(t-1)/2\), and
\[
f_t(u)=L_0\prod_{j=1}^t\frac{u^2+1}{u^2+j}.
\]
This is nondecreasing for \(u\ge0\), lies in \((0,L_0)\), and
\[
f_t(u)\ge L_0\left(1-\frac{a}{u^2+1}\right).
\]
The half-width condition excludes \(F\le0\) and \(F\ge L_0+1\).
For \(1\le F\le t-1\),
\(f_t(F)\ge2^t/(t+1)>t-1/2\ge F+1/2\).
For \(t\le F\le L_0/2\), the displayed lower bound exceeds
\(L_0/2+1/2\). For \(L_0/2\le F\le L_0-1\), it gives
\(L_0-f_t(F)<4a/L_0=2/(t-2)!<1/2\), again excluding \(F\).
At \(F=L_0\),
\[
0<L_0-f_t(L_0)<\frac{a}{L_0}=
\frac1{2(t-2)!}<\frac1{32},
\]
which proves both the characterization and the completeness margin.

We have now recovered \(F=t!\), \(B>0\), \(q>0\), and \(y=V-1>F+2\).
The endpoint identities are
\[
x=F+1+(y+1)u_0,\qquad
x=1+\bigl(1+(t+1)y\bigr)u_t.
\]
Since \(x\ge t>1\), the second gives \(u_t\ge1\), hence
\begin{equation}\label{eq:endpoint}
x\ge2+(t+1)y.
\end{equation}
The inherited coefficient estimates recover all remaining digit ranges.
The lowered positions in Section~\ref{sec:strict} only strengthen them.
The strict amplification proof, using \(R=2+t+x+y+F<2x\), the bound
\(x>2^{32}\), and the retained majorant exponent \(2^{23}\), again gives
\(R^{2^{23}}<2^x\). The finite CRT and Horn-flow decoding therefore
produce a C222 refutation.

\section{Completeness from a finite refutation}
\label{sec:completeness}

Start with a finite C222 refutation. Select the strict cofinal horizon as
in Section~\ref{sec:strict}; then define \(F=t!\), \(V=(F^2+1)^t\),
and \(y=V-1\). Section~\ref{sec:beta} supplies the pending-state bound
and beta CRT hypotheses. The parent's construction gives the beta code
\(x\), the exact factorial-binomial value \(B\), and the nine CRT
representatives at precision \(N^2\), where \(N=B^x\).

Choose representatives cofinally in their residue classes, as in the
retained semantic-anchor construction: the dominant schema representative
can be between \(2N^2\) and \(4N^2\), with prescribed parity, while the
others are at most \(2N^2\). This makes the large global derived index
\(h\) even and strictly larger than every other principal index. This
choice does not alter the small local values or their strict bounds.
Set \(H=1+2\sum a_i\), \(q=(BH)^x\), and fill all seventeen genuine
power values.

Only now form \(W\) from the semantic atoms. The ordinary parent frame
chooses a sufficiently large canonical Pell abscissa at center \(W+1\),
with ordinate an odd multiple of \(W\). It gives \(|A|>W^W\).
The canonical principal pairs, squared semantic anchor, common auxiliary
pair, and comparison pair exist by the retained constructive Pell and
finite-period arguments. No new condition on the frame is imposed.

Every exact power satisfies its original power congruence, so
Lemma~\ref{lem:identity} supplies the two aggregate integers for each
group. Thus completeness of the new pooled rows follows from simultaneous
completeness of the old congruences without requiring that their
individual quotient coordinates remain in the equation.

Every true digit remainder, normalized to \(M=q^x\), is smaller than
\(M/32\); the rounded factorial error has the same bound. The twelve
errors therefore satisfy
\[
4\left(e^2+\sum_{j=1}^{11}\rho_j^2\right)
 <\frac{48}{1024}M^2<M^2.
\]
The remaining nonnegative integer slack is represented by
\(\G(u,v,w)=u^2+v^2+w(w+1)\). One quick proof of its universality is
to apply the three-square theorem to \(4n+1\): exactly one of the
three square roots is odd, giving the claimed form after division by four.
The other retained guard is filled in the same way.

Finally apply the exact coordinate maps of the next section. Every finite
refutation thus yields an integer zero of the delivered polynomial.
Together with soundness this proves the required equivalence of
existential statements. The strict-horizon and beta changes go through
freshly constructed refutation witnesses; they are not claimed to be
bijections between arbitrary old and new assignments at a fixed horizon.

\section{Exact reductions of auxiliary coordinates}
\label{sec:algebra}

Three further variables disappear by identities over \(\Z\). These
steps have explicit maps in both directions and do not need semantic
bounds, positivity, or divisions.

\subsection{Store a sum when only the sum is used}

Suppose \(u=f(\mathbf z)\), \(v=g(\mathbf z)\) are private definitions,
and every other residual uses these variables only through \(u+v\).
Replace them by \(w=f(\mathbf z)+g(\mathbf z)\). Projection gives
\(w=u+v\), and restoration gives \(u=f\), \(v=g\). If \(f,g\)
are quadratic, their sum is still quadratic even when substituting either
separately into a product would have raised degree.

In the centered parent, the two values are
\[
x_{81}=x_{80}^2,\qquad
x_{83}=(x_{14}-x_{15}-x_{82})^2.
\]
Their only external use is
\(-x_{101}(x_{81}+x_{83})\). One sum coordinate therefore replaces
the two and removes one variable and one row.

\subsection{Compute the products that the remaining rows use}

Let \(F=x_3\), \(S=x_{14}\), \(S'=x_{15}\). The product
\(FS'=x_{70}\) already exists. Three auxiliary definitions in the
centered parent are
\[
x_{84}=1+Fx_{11},\quad x_{85}=1+F(x_{10}-1),\quad
x_{89}=1+F(6x_{10}+8x_{11}+2x_{21}).
\]
They are used only after multiplication by \(S\) or \(S'\).
Introduce one coordinate \(J=FS\), delete these three definitions,
and use the distributive identities
\begin{align*}
S'(1+Fx_{11})&=S'+(FS')x_{11},\\
S(1+F(x_{10}-1))&=S+J(x_{10}-1),\\
S(1+Fx_{11})&=S+Jx_{11},\\
S'(1+F\ell)&=S'+(FS')\ell,
\qquad\ell=6x_{10}+8x_{11}+2x_{21}.
\end{align*}
All replacement residuals remain quadratic. The old values restore by
their three definitions, so this saves two variables and two rows with
an exact integral inverse on zero sets.

\subsection{Pell difference coordinates reduce density}

To control the density caused by pooling, replace each principal
abscissa by the difference
\(D_i=X_i-Z_i\). This is a triangular integral change of coordinates,
with inverse \(X_i=D_i+Z_i\). The norm becomes
\[
D_i^2+2D_iZ_i+Y_i^2-1=0,
\]
and the long pooled square uses \((m_i-D_i)^2\).
Before that substitution, reduce the orientation row by the retained
norms. Modulo those norms its summand satisfies
\[
(X_i-1)^2+(Y_i-k_i)^2
\equiv Z_i^2+k_i^2-2X_i-2k_iY_i+2.
\]
After the difference substitution, \(D_i\) occurs only linearly here.
Adding or subtracting a retained residual from another is reversible
and preserves the common zero set. The independent checker authenticates
the exact row combinations and every offset in the centered coordinates.

These transformations yield the following total saving:
\[
\underbrace{4}_{\text{strict ranges}}+
\underbrace{1}_{\text{beta base}}+
\underbrace{3}_{\text{auxiliary coordinates}}+
\underbrace{12}_{\text{power groups}}=20.
\]
The last density change does not alter the number of variables.

\section{Calibration and the formalization boundary}
\label{sec:calibration}

The underlying source is the numerical proof calculus C222, under the
\texttt{xx0} square-shell recoding used by the direct arithmetic
construction. The
constructors, the eight ordinary axiom families, the guarded bundle,
modus ponens, and generalization are retained exactly from the immediate
parent. Falsity has code 1 and truth code 0. The new equation
changes how finite refutation certificates are represented arithmetically,
not which calculus is calibrated.

The semantic chain is
\[
\begin{split}
\exists\mathbf x\ P_{\NumVars}(\mathbf x)=0
&\quad\longleftrightarrow\quad
\text{a finite square-shell C222 derivation of 1}\\
&\quad\longleftrightarrow\quad
\text{a finite original C222 derivation of 1}\\
&\quad\longleftrightarrow\quad
\text{a finite ZF refutation}\\
&\quad\longleftrightarrow\quad
\text{a finite ZFC refutation}.
\end{split}
\]

For clarity, the recoding is an ordinary theorem rather than a claim that
the two sets of constructor polynomials are identical. Put
\[
Q_0(a,b)=(a+b)^2+a,\qquad
Q_1(a,b)=(a+b)^2+3a+4b+3.
\]
The shell constructors are \(I=2Q_0\), \(M=2Q_0+1\),
\(\mathsf A=2Q_1\), and \(\mathsf E=2Q_1+1\).
On the shell \(a+b=s\), \(Q_0\) covers
\([s^2,s^2+s]\) once; on the preceding shell, \(Q_1\) covers
\([s^2+s+1,(s+1)^2-1]\) once. These intervals partition \(\N\).
Every positive constructor value exceeds both arguments. Decoding and
recursing therefore gives a total constructor-preserving bijection
\(f\) from the original Cantor coding to this shell coding, with the
shared case \(I(0,0)=0\) treated as a base case. Its inverse is obtained
by the opposite decoding; induction proves both inverse identities.
In particular \(f(0)=0\), \(f(1)=1\), and equality and disequality are
preserved.

All constructor-term axioms and the two inference rules transport by
\(f\). In the guarded branch, the arithmetic witness must be rebuilt:
from the two transported freshness disequalities choose
\[
d_{\mathrm{new}}
=\left|(f(c)-f(a))(f(c)-f(b))\right|-1\ge0.
\]
This gives the required squared guard. It does not assume that \(f\)
preserves addition. A separate branch in which \(d\) is a code parameter
transports \(d\) normally; the branches are alternatives, so these
assignments do not conflict. The inverse bijection proves the converse
transport. The full inherited argument is supplied unmodified in
\path{inherited/shell_recoding_supplement/SHELL_ENCODING_PROOF.md}, with
the literal constructor sources and their pinned provenance.

The named logical interfaces in the retained source are
\path{KEndpoint.c222_iff_zf_refutation} and
\path{ConstructibleChoice.zf_refutation_iff_zfc_refutation}.
The latter uses the constructible-universe equiconsistency argument.
The relevant interface files, source pins, ordinary arithmetic proofs,
and original notices are included under
\path{inherited/joint_norm245/}; the separately identified syntax supplement
closes the recoding step. These retained sources are not a complete
recompiled proof-assistant dependency closure.

For auditing the new arithmetic, the most relevant retained proof files
are the Joint Norm \path{proofs/BOOTSTRAP_PROOF.md} and, within its
normalized-parent closure, \path{PROOF_315.md},
\path{ORIENTATION_FIRST_PROOF.md}, \path{STRICT_BUDGET_PROOF.md},
\path{INHERITED_DERIVED_LABEL_RANGE_PROOF.md}, and the direct calculus
and CRT exposition. The provenance manifest gives exact file paths and
hashes, rather than relying on these abbreviated names to identify a
dependency.

All new arguments use integer polynomial identities, induction on finite
recurrences and proofs, elementary bounds, finite CRT, and finite-group
Pell periodicity. They are ordinary mathematics formalizable in ZFC.
Together with the identified inherited calibration they prove
Theorem~\ref{thm:main}. The supplied Python checks certify finite objects
and exact transformations. They do not certify the absence of zeros,
establish consistency, or replace the unbounded mathematical proofs.
The full imported Lean dependency closure was not recompiled in this
work, and no new end-to-end Lean or Rocq theorem is claimed.

\section{Reproduction and finite evidence}

The package contains the literal residual system, its full sparse
expansion, the construction stages, transformation certificates, the
unmodified immediate parent and retained proof closure, the report source,
and independent verification scripts. Python 3 with its standard library
is sufficient for the new finite checks.

From the package root, run
\begin{verbatim}
python3 -B verify.py
python3 -B -O verify.py
\end{verbatim}
The build command and the PDF rebuild command are in \path{README.md}.
The new verifier checks source identities, the semantic edit boundaries,
the auxiliary-coordinate maps, the two power groups and their exact
affine ports, the unchanged square frame and atom coverage, the Pell
coordinate changes, and every coefficient of the collected polynomial.
It computes the reported dimensions, degree, support, coefficient height,
and coordinate liveness from the delivered residuals.

Separate exact checks exercise the aggregate polynomial identity, signed
Pell congruences, unit and negative bases, negative indices, the square-frame
error estimate, and the corrected zero-factorial bound. These are useful
falsification and regression checks; the preceding theorems establish the
unbounded conclusions. The reviews were independently assigned to other
AI research agents, not represented as external human peer review.

\section{Further research questions}

\begin{enumerate}
\item \textbf{Can distinct bases share one aggregate?}
The present proof groups by an identical base and uses its single modulus
\(2An-n^2-1\). Pooling different bases without new product or CRT
coordinates would require a different quadratic identity. A naive sum
of their errors has no common modulus and is not justified by this proof.
\item \textbf{Can an entire Pell pair disappear?}
We removed quotient coordinates while keeping canonical pairs. A useful
next target is a semantic exponent identity that reconstructs a pair
already needed elsewhere, with no new multiplication coordinates.
\item \textbf{How far can auxiliary-coordinate compression go?}
The sum replacement and product redistribution exploit how values are
used by the remaining equations. A systematic search over integral
changes of coordinates and quadratic row combinations may find more
reductions than single-row monic elimination. Coefficient height and
support must be measured alongside dimension.
\item \textbf{Can the new module be formally verified independently?}
The aggregate identity, signed modular power lemma, and square-frame
error bound form a small coherent target for Lean or Rocq. A theorem
with an explicit list of affine ports would allow subsequent compilations
to reuse one checked module.
\item \textbf{What is the best calculus for the arithmetic cost?}
Optimizing the number of schema parameters can help, but a parameter
hidden in a marker premise may reappear as an additional encoded column
in the global proof-flow construction. A local reduction must be carried
through the complete arithmetic compiler before claiming a smaller
consistency equation. The full guarded-bundle freshness range must also
be retained unless a new logical calibration is proved.
\item \textbf{Can density be reduced without giving back the variables?}
The principal equation has more collected monomials than its parent.
Alternative quadratic residual bases and integral Pell coordinates are
natural targets for a future support/height optimization at dimension
\NumVars.
\end{enumerate}

\section*{Attribution}
\addcontentsline{toc}{section}{Attribution}
The C222 calibration, finite flow/CRT encoding, rounded factorial, signed
Pell cap, cofinal semantic anchor, and common digit norm are inherited
and identified in the source ledger. The new contributions are the
two-equation pooling of exact powers at a common base with no modulus
coordinate, its signed square-frame proof and integration into the
bootstrap, the strict-horizon and linear beta refinements, the exact
auxiliary-coordinate reductions, and the literal combined endpoint.
No claim is made that elementary polynomial identities themselves are
new. Novelty is stated relative to the inspected constructions; a global
priority claim would require a broader literature review.

\begin{thebibliography}{9}
\small\setlength{\itemsep}{2pt}
\bibitem{repo} V. Reshetnikov, \emph{ProveIt}, snapshot
\texttt{a0ef2df5bb1f73c35cd334bbc7f53553350ef06a}, 25 October 2026.
\url{https://github.com/VladimirReshetnikov/ProveIt/commit/a0ef2df5bb1f73c35cd334bbc7f53553350ef06a}.
\bibitem{joint} \emph{A single norm for factorial and digit certificates},
incoming package \path{ConZFC_Joint_Norm_2026-10-25.zip}, in the preceding
repository snapshot. The unmodified direct component is included here.
\url{https://github.com/VladimirReshetnikov/ProveIt/blob/a0ef2df5bb1f73c35cd334bbc7f53553350ef06a/docs/incoming/ConZFC_Joint_Norm_2026-10-25.zip}.
\bibitem{forms} V. Reshetnikov, \emph{Con(ZFC) in Diophantine form},
collection README and placement ledger, same snapshot.
\url{https://github.com/VladimirReshetnikov/ProveIt/tree/a0ef2df5bb1f73c35cd334bbc7f53553350ef06a/Computability/DiophantineConsistency/con-zfc-diophantine-forms}.
\bibitem{batched} \emph{Shared arithmetic certificates for Diophantine
consistency equations}, collection package 96, especially
\path{report/batched_powers.tex}. This earlier batching shares the Pell
environment and retains separate local power congruences.
\bibitem{pak} K. P\k{a}k and C. Kaliszyk, \emph{Formalizing a Diophantine
Representation of the Set of Prime Numbers}, ITP 2022,
arXiv:2204.12311. \url{https://arxiv.org/abs/2204.12311}.
\bibitem{sun} Z.-W. Sun, \emph{Ten unknowns for Hilbert's tenth problem
over the integers}, arXiv:2609.01594v2, 2026.
\url{https://arxiv.org/abs/2609.01594v2}.
\end{thebibliography}

\appendix
\clearpage
\section{Literal residual specification}\label{app:residuals}
The following \NumRows\ residuals, in precisely this order, define
\eqref{eq:main}. All \(x_i\) range independently over \(\Z\). A line
break within a residual does not start a new equation. This appendix is
generated directly from \path{equation/equation.json} and checked against
that file. The variable names in the machine-readable list are provenance
labels; only the displayed integer coefficients and indices define the
polynomial. Every index from 0 through \LastVar\ occurs.
% Generated directly from equation/equation.json.
\begingroup
\small
\interdisplaylinepenalty=10000
\setlength{\jot}{3pt}
% residual 0
\begin{align*}
R_{0}={}& -x_{2}^{2} + x_{3}
\end{align*}
% residual 1
\begin{align*}
R_{1}={}& x_{0} - x_{0}x_{5} - x_{5} + x_{8}
\end{align*}
% residual 2
\begin{align*}
R_{2}={}& -1 + x_{1} - x_{2} - x_{5}x_{6}
\end{align*}
% residual 3
\begin{align*}
R_{3}={}& -1 + x_{1} - x_{7}x_{8}
\end{align*}
% residual 4
\begin{align*}
R_{4}={}& 1 + x_{3} - x_{4}x_{19} + x_{17}
\end{align*}
% residual 5
\begin{align*}
R_{5}={}& -1 + 2x_{9} - x_{9}^{2} - 2x_{9}x_{10} + 2x_{10} - x_{10}^{2} + x_{20}
\end{align*}
% residual 6
\begin{align*}
R_{6}={}& -5 - 30x_{9} - 36x_{9}^{2} - 96x_{9}x_{10} - 24x_{9}x_{20} - 40x_{10} - 64x_{10}^{2} - 32x_{10}x_{20} - 10x_{20} \\
&{} - 4x_{20}^{2} + x_{21}
\end{align*}
% residual 7
\begin{align*}
R_{7}={}& -2 + 2x_{9} - x_{9}^{2} - 4x_{9}x_{21} + 4x_{21} - 4x_{21}^{2} + x_{22}
\end{align*}
% residual 8
\begin{align*}
R_{8}={}& 1 + 6x_{9} - 36x_{9}^{2} - 192x_{9}x_{21} - 24x_{9}x_{22} + 16x_{21} - 256x_{21}^{2} - 64x_{21}x_{22} + 2x_{22} \\
&{} - 4x_{22}^{2} + x_{23}
\end{align*}
% residual 9
\begin{align*}
R_{9}={}& -2 - x_{10}^{2} - 2x_{10}x_{11} - x_{11}^{2} + x_{24}
\end{align*}
% residual 10
\begin{align*}
R_{10}={}& -9 + 7x_{9} - x_{9}^{2} - 4x_{9}x_{10} - 4x_{9}x_{24} + 12x_{10} - 4x_{10}^{2} - 8x_{10}x_{24} + 12x_{24} - 4x_{24}^{2} \\
&{} + x_{25}
\end{align*}
% residual 11
\begin{align*}
R_{11}={}& -1 + 2x_{9} - x_{9}^{2} - 2x_{9}x_{11} + 2x_{11} - x_{11}^{2} + x_{26}
\end{align*}
% residual 12
\begin{align*}
R_{12}={}& -6x_{9} - 144x_{9}^{2} - 192x_{9}x_{10} - 192x_{9}x_{11} - 48x_{9}x_{20} - 48x_{9}x_{26} - 8x_{10} - 64x_{10}^{2} \\
&{} - 128x_{10}x_{11} - 32x_{10}x_{20} - 32x_{10}x_{26} - 64x_{11}^{2} - 32x_{11}x_{20} - 32x_{11}x_{26} - 2x_{20} \\
&{} - 4x_{20}^{2} - 8x_{20}x_{26} - 4x_{26}^{2} + x_{27}
\end{align*}
% residual 13
\begin{align*}
R_{13}={}& -4x_{25}^{2} - 8x_{25}x_{27} - 4x_{27}^{2} + x_{28}
\end{align*}
% residual 14
\begin{align*}
R_{14}={}& -6 + 16x_{11} - 16x_{11}^{2} - 16x_{11}x_{24} - 16x_{11}x_{26} + 8x_{24} - 4x_{24}^{2} - 8x_{24}x_{26} + 8x_{26} \\
&{} - 4x_{26}^{2} + x_{29}
\end{align*}
% residual 15
\begin{align*}
R_{15}={}& 1 + 6x_{9} - 36x_{9}^{2} - 96x_{9}x_{10} - 48x_{9}x_{11} - 24x_{9}x_{20} - 48x_{9}x_{26} - 24x_{9}x_{29} + 8x_{10} \\
&{} - 64x_{10}^{2} - 64x_{10}x_{11} - 32x_{10}x_{20} - 64x_{10}x_{26} - 32x_{10}x_{29} + 8x_{11} - 16x_{11}^{2} \\
&{} - 16x_{11}x_{20} - 32x_{11}x_{26} - 16x_{11}x_{29} + 2x_{20} - 4x_{20}^{2} - 16x_{20}x_{26} - 8x_{20}x_{29} + 8x_{26} \\
&{} - 16x_{26}^{2} - 16x_{26}x_{29} + 4x_{29} - 4x_{29}^{2} + x_{30}
\end{align*}
% residual 16
\begin{align*}
R_{16}={}& -10 - 5x_{9} - x_{9}^{2} - 12x_{9}x_{10} - 16x_{9}x_{11} - 4x_{9}x_{24} - 36x_{10} - 36x_{10}^{2} - 96x_{10}x_{11} \\
&{} - 24x_{10}x_{24} - 48x_{11} - 64x_{11}^{2} - 32x_{11}x_{24} - 12x_{24} - 4x_{24}^{2} + x_{31}
\end{align*}
% residual 17
\begin{align*}
R_{17}={}& -2 - 24x_{9} - 36x_{9}^{2} - 12x_{9}x_{10} - 96x_{9}x_{11} - 24x_{9}x_{26} - 3x_{10} - x_{10}^{2} - 16x_{10}x_{11} \\
&{} - 4x_{10}x_{26} - 32x_{11} - 64x_{11}^{2} - 32x_{11}x_{26} - 8x_{26} - 4x_{26}^{2} + x_{32}
\end{align*}
% residual 18
\begin{align*}
R_{18}={}& 2 - 4x_{31}^{2} - 8x_{31}x_{32} - 4x_{32}^{2} + x_{33}
\end{align*}
% residual 19
\begin{align*}
R_{19}={}& -12 + 16x_{9} - 4x_{9}^{2} - 8x_{9}x_{11} - 8x_{9}x_{20} - 8x_{9}x_{24} + 16x_{11} - 4x_{11}^{2} - 8x_{11}x_{20} \\
&{} - 8x_{11}x_{24} + 16x_{20} - 4x_{20}^{2} - 8x_{20}x_{24} + 16x_{24} - 4x_{24}^{2} + x_{34}
\end{align*}
% residual 20
\begin{align*}
R_{20}={}& -3 - 24x_{9} - 64x_{9}^{2} - 160x_{9}x_{11} - 64x_{9}x_{24} - 32x_{9}x_{26} - 32x_{9}x_{34} - 26x_{11} - 100x_{11}^{2} \\
&{} - 80x_{11}x_{24} - 40x_{11}x_{26} - 40x_{11}x_{34} - 8x_{24} - 16x_{24}^{2} - 16x_{24}x_{26} - 16x_{24}x_{34} - 6x_{26} \\
&{} - 4x_{26}^{2} - 8x_{26}x_{34} - 4x_{34} - 4x_{34}^{2} + x_{35}
\end{align*}
% residual 21
\begin{align*}
R_{21}={}& -1 - 18x_{9} - 36x_{9}^{2} - 96x_{9}x_{10} - 24x_{9}x_{20} - 24x_{10} - 64x_{10}^{2} - 32x_{10}x_{20} - 6x_{20} \\
&{} - 4x_{20}^{2} + x_{36}
\end{align*}
% residual 22
\begin{align*}
R_{22}={}& -2 + 2x_{9} - x_{9}^{2} - 4x_{9}x_{36} + 4x_{36} - 4x_{36}^{2} + x_{37}
\end{align*}
% residual 23
\begin{align*}
R_{23}={}& 1 + 6x_{9} - 36x_{9}^{2} - 192x_{9}x_{36} - 24x_{9}x_{37} + 16x_{36} - 256x_{36}^{2} - 64x_{36}x_{37} + 2x_{37} \\
&{} - 4x_{37}^{2} + x_{38}
\end{align*}
% residual 24
\begin{align*}
R_{24}={}& -x_{10}^{2} - 4x_{10}x_{38} - 4x_{38}^{2} + x_{39}
\end{align*}
% residual 25
\begin{align*}
R_{25}={}& -6 - 30x_{9} - 64x_{9}^{2} - 128x_{9}x_{10} - 64x_{9}x_{20} - 32x_{10} - 64x_{10}^{2} - 64x_{10}x_{20} - 14x_{20} \\
&{} - 16x_{20}^{2} + x_{40}
\end{align*}
% residual 26
\begin{align*}
R_{26}={}& -4 + 8x_{10} - 64x_{10}^{2} - 96x_{10}x_{11} - 32x_{10}x_{24} - 32x_{10}x_{40} + 6x_{11} - 36x_{11}^{2} \\
&{} - 24x_{11}x_{24} - 24x_{11}x_{40} + 2x_{24} - 4x_{24}^{2} - 8x_{24}x_{40} + 4x_{40} - 4x_{40}^{2} + x_{41}
\end{align*}
% residual 27
\begin{align*}
R_{27}={}& 8x_{9} - 64x_{9}^{2} - 96x_{9}x_{11} - 32x_{9}x_{26} - 32x_{9}x_{41} + 6x_{11} - 36x_{11}^{2} - 24x_{11}x_{26} \\
&{} - 24x_{11}x_{41} + 2x_{26} - 4x_{26}^{2} - 8x_{26}x_{41} + 4x_{41} - 4x_{41}^{2} + x_{42}
\end{align*}
% residual 28
\begin{align*}
R_{28}={}& -x_{11}^{2} - 2x_{11}x_{12} - x_{12}^{2} + x_{43}
\end{align*}
% residual 29
\begin{align*}
R_{29}={}& 2x_{9} - x_{9}^{2} - 4x_{9}x_{11} - 4x_{9}x_{43} + 4x_{11} - 4x_{11}^{2} - 8x_{11}x_{43} + 4x_{43} - 4x_{43}^{2} + x_{44}
\end{align*}
% residual 30
\begin{align*}
R_{30}={}& -35 + 24x_{9} - 4x_{9}^{2} - 8x_{9}x_{10} - 8x_{9}x_{24} - 8x_{9}x_{26} + 24x_{10} - 4x_{10}^{2} - 8x_{10}x_{24} \\
&{} - 8x_{10}x_{26} + 24x_{24} - 4x_{24}^{2} - 8x_{24}x_{26} + 24x_{26} - 4x_{26}^{2} + x_{45}
\end{align*}
% residual 31
\begin{align*}
R_{31}={}& -16x_{9}^{2} - 32x_{9}x_{11} - 32x_{9}x_{26} - 32x_{9}x_{43} - 16x_{9}x_{44} - 16x_{9}x_{45} - 4x_{11} - 16x_{11}^{2} \\
&{} - 32x_{11}x_{26} - 32x_{11}x_{43} - 16x_{11}x_{44} - 16x_{11}x_{45} - 16x_{26}^{2} - 32x_{26}x_{43} - 16x_{26}x_{44} \\
&{} - 16x_{26}x_{45} - 4x_{43} - 16x_{43}^{2} - 16x_{43}x_{44} - 16x_{43}x_{45} - 2x_{44} - 4x_{44}^{2} - 8x_{44}x_{45} \\
&{} - 4x_{45}^{2} + x_{46}
\end{align*}
% residual 32
\begin{align*}
R_{32}={}& x_{11} - 25x_{11}^{2} - 40x_{11}x_{26} - 20x_{11}x_{29} - 16x_{26}^{2} - 16x_{26}x_{29} - 4x_{29}^{2} + x_{47}
\end{align*}
% residual 33
\begin{align*}
R_{33}={}& -1 + 8x_{10} - 4x_{10}^{2} - 8x_{10}x_{24} + 8x_{24} - 4x_{24}^{2} + x_{48}
\end{align*}
% residual 34
\begin{align*}
R_{34}={}& -1 - 16x_{10}^{2} - 8x_{10}x_{11} - 32x_{10}x_{24} - 16x_{10}x_{48} - x_{11}^{2} - 8x_{11}x_{24} - 4x_{11}x_{48} \\
&{} - 16x_{24}^{2} - 16x_{24}x_{48} - 4x_{48}^{2} + x_{49}
\end{align*}
% residual 35
\begin{align*}
R_{35}={}& 1 + x_{10} - x_{10}^{2} - 4x_{10}x_{11} - 4x_{10}x_{49} - 4x_{11}^{2} - 8x_{11}x_{49} - 4x_{49}^{2} + x_{50}
\end{align*}
% residual 36
\begin{align*}
R_{36}={}& -1 - 4x_{50} - 4x_{50}^{2} + x_{51}
\end{align*}
% residual 37
\begin{align*}
R_{37}={}& -1 - x_{11}^{2} - 8x_{11}x_{50} - 4x_{11}x_{51} - 16x_{50}^{2} - 16x_{50}x_{51} - 4x_{51}^{2} + x_{52}
\end{align*}
% residual 38
\begin{align*}
R_{38}={}& 2 + x_{10} - x_{10}^{2} - 4x_{10}x_{11} - 4x_{10}x_{52} - 4x_{11}^{2} - 8x_{11}x_{52} - 4x_{52}^{2} + x_{53}
\end{align*}
% residual 39
\begin{align*}
R_{39}={}& -4x_{47}^{2} - 8x_{47}x_{53} - 4x_{53}^{2} + x_{54}
\end{align*}
% residual 40
\begin{align*}
R_{40}={}& -9 + 7x_{9} - x_{9}^{2} - 8x_{9}x_{47} - 4x_{9}x_{54} + 24x_{47} - 16x_{47}^{2} - 16x_{47}x_{54} + 12x_{54} - 4x_{54}^{2} \\
&{} + x_{55}
\end{align*}
% residual 41
\begin{align*}
R_{41}={}& -1 - 4x_{55} - 4x_{55}^{2} + x_{56}
\end{align*}
% residual 42
\begin{align*}
R_{42}={}& -x_{11}^{2} - 8x_{11}x_{55} - 4x_{11}x_{56} - 16x_{55}^{2} - 16x_{55}x_{56} - 4x_{56}^{2} + x_{57}
\end{align*}
% residual 43
\begin{align*}
R_{43}={}& -x_{9} - x_{9}^{2} - 4x_{9}x_{11} - 4x_{9}x_{57} - 4x_{11} - 4x_{11}^{2} - 8x_{11}x_{57} - 4x_{57} - 4x_{57}^{2} + x_{58}
\end{align*}
% residual 44
\begin{align*}
R_{44}={}& -10 - 12x_{58} - 4x_{58}^{2} + x_{59}
\end{align*}
% residual 45
\begin{align*}
R_{45}={}& 1 + 4x_{10} - 4x_{10}^{2} - 8x_{10}x_{24} - 16x_{10}x_{58} - 8x_{10}x_{59} + 4x_{24} - 4x_{24}^{2} - 16x_{24}x_{58} \\
&{} - 8x_{24}x_{59} + 8x_{58} - 16x_{58}^{2} - 16x_{58}x_{59} + 4x_{59} - 4x_{59}^{2} + x_{60}
\end{align*}
% residual 46
\begin{align*}
R_{46}={}& x_{10} - 25x_{10}^{2} - 40x_{10}x_{24} - 20x_{10}x_{60} - 16x_{24}^{2} - 16x_{24}x_{60} - 4x_{60}^{2} + x_{61}
\end{align*}
% residual 47
\begin{align*}
R_{47}={}& -10 - 12x_{53} - 4x_{53}^{2} + x_{62}
\end{align*}
% residual 48
\begin{align*}
R_{48}={}& 1 - 16x_{53}^{2} - 16x_{53}x_{61} - 16x_{53}x_{62} - 4x_{61}^{2} - 8x_{61}x_{62} - 4x_{62}^{2} + x_{63}
\end{align*}
% residual 49
\begin{align*}
R_{49}={}& x_{11} - x_{11}^{2} - 8x_{11}x_{61} - 4x_{11}x_{63} - 16x_{61}^{2} - 16x_{61}x_{63} - 4x_{63}^{2} + x_{64}
\end{align*}
% residual 50
\begin{align*}
R_{50}={}& -4x_{9} - 4x_{9}^{2} - 4x_{9}x_{11} - 8x_{9}x_{20} - x_{11} - x_{11}^{2} - 4x_{11}x_{20} - 4x_{20} - 4x_{20}^{2} + x_{65}
\end{align*}
% residual 51
\begin{align*}
R_{51}={}& 1 + 4x_{9} - 4x_{9}^{2} - 8x_{9}x_{20} - 8x_{9}x_{65} + 4x_{20} - 4x_{20}^{2} - 8x_{20}x_{65} + 4x_{65} - 4x_{65}^{2} \\
&{} + x_{66}
\end{align*}
% residual 52
\begin{align*}
R_{52}={}& -2 - 4x_{9} - 16x_{9}^{2} - 32x_{9}x_{20} - 16x_{9}x_{64} - 16x_{9}x_{66} - 4x_{20} - 16x_{20}^{2} - 16x_{20}x_{64} \\
&{} - 16x_{20}x_{66} - 4x_{64}^{2} - 8x_{64}x_{66} - 2x_{66} - 4x_{66}^{2} + x_{67}
\end{align*}
% residual 53
\begin{align*}
R_{53}={}& -2 - 4x_{67} - 4x_{67}^{2} + x_{68}
\end{align*}
% residual 54
\begin{align*}
R_{54}={}& -x_{2}x_{14} + x_{69}
\end{align*}
% residual 55
\begin{align*}
R_{55}={}& x_{13} - x_{14} - 2x_{23}x_{69} - x_{70}
\end{align*}
% residual 56
\begin{align*}
R_{56}={}& x_{13} - x_{14} - 4x_{25}x_{69} - 2x_{28}x_{69} - x_{71}
\end{align*}
% residual 57
\begin{align*}
R_{57}={}& x_{13} - x_{14} - 2x_{30}x_{69} - 6x_{69} - x_{72}
\end{align*}
% residual 58
\begin{align*}
R_{58}={}& x_{13} - x_{14} - 4x_{31}x_{69} - 2x_{33}x_{69} - x_{73}
\end{align*}
% residual 59
\begin{align*}
R_{59}={}& x_{13} - x_{14} - 2x_{35}x_{69} + 6x_{69} - x_{74}
\end{align*}
% residual 60
\begin{align*}
R_{60}={}& x_{13} - x_{14} - 4x_{38}x_{69} - 2x_{39}x_{69} - x_{75}
\end{align*}
% residual 61
\begin{align*}
R_{61}={}& x_{13} - x_{14} - 2x_{42}x_{69} - x_{76}
\end{align*}
% residual 62
\begin{align*}
R_{62}={}& x_{13} - x_{14} - 2x_{46}x_{69} - 4x_{69} - x_{77}
\end{align*}
% residual 63
\begin{align*}
R_{63}={}& -x_{9}x_{10} + x_{9}x_{11} + x_{10} + x_{10}x_{11} - x_{11} - x_{11}^{2} + x_{78}
\end{align*}
% residual 64
\begin{align*}
R_{64}={}& 1 + 2x_{12} + x_{12}^{2} - x_{78}^{2} + x_{79}
\end{align*}
% residual 65
\begin{align*}
R_{65}={}& -x_{13}^{2} + 2x_{13}x_{14} + 2x_{13}x_{81} - x_{14}^{2} - 2x_{14}x_{81} - x_{79}^{2} + x_{80} - x_{81}^{2}
\end{align*}
% residual 66
\begin{align*}
R_{66}={}& -4x_{67}x_{69} - 2x_{68}x_{69} + 2x_{69} + x_{81}
\end{align*}
% residual 67
\begin{align*}
R_{67}={}& -x_{10}x_{69} - x_{14} + x_{82}x_{83} + x_{84}
\end{align*}
% residual 68
\begin{align*}
R_{68}={}& -1 + 2x_{2} - 2x_{2}x_{9} - 2x_{2}x_{20} + x_{82}
\end{align*}
% residual 69
\begin{align*}
R_{69}={}& -x_{9}x_{220} - x_{13} + x_{83} + x_{220}
\end{align*}
% residual 70
\begin{align*}
R_{70}={}& -6x_{9}x_{69} - 8x_{10}x_{69} + x_{10}x_{220} + x_{13} - x_{14} - 2x_{20}x_{69} + x_{85}
\end{align*}
% residual 71
\begin{align*}
R_{71}={}& -x_{13}x_{70} + x_{14}x_{70} + x_{86}
\end{align*}
% residual 72
\begin{align*}
R_{72}={}& -x_{71}x_{86} + x_{87}
\end{align*}
% residual 73
\begin{align*}
R_{73}={}& -x_{72}x_{87} + x_{88}
\end{align*}
% residual 74
\begin{align*}
R_{74}={}& -x_{73}x_{88} + x_{89}
\end{align*}
% residual 75
\begin{align*}
R_{75}={}& -x_{74}x_{89} + x_{90}
\end{align*}
% residual 76
\begin{align*}
R_{76}={}& -x_{75}x_{90} + x_{91}
\end{align*}
% residual 77
\begin{align*}
R_{77}={}& -x_{76}x_{91} + x_{92}
\end{align*}
% residual 78
\begin{align*}
R_{78}={}& -x_{77}x_{92} + x_{93}
\end{align*}
% residual 79
\begin{align*}
R_{79}={}& -x_{84}x_{93} + x_{94}
\end{align*}
% residual 80
\begin{align*}
R_{80}={}& -x_{85}x_{94} + x_{95}
\end{align*}
% residual 81
\begin{align*}
R_{81}={}& -x_{95}^{2} + x_{96}
\end{align*}
% residual 82
\begin{align*}
R_{82}={}& -x_{5} - x_{5}x_{17} + x_{17} + x_{97}
\end{align*}
% residual 83
\begin{align*}
R_{83}={}& -x_{15}x_{97} + x_{98}
\end{align*}
% residual 84
\begin{align*}
R_{84}={}& -2x_{1}^{2} + 2x_{1}x_{13} + 2x_{1}x_{14} + 2x_{1}x_{98} + 2x_{1}x_{99} - x_{13}^{2} - 2x_{13}x_{98} - x_{14}^{2} - 2x_{14}x_{99} \\
&{} + x_{18}x_{101} - x_{80}x_{96} - x_{98}^{2} - x_{99}^{2}
\end{align*}
% residual 85
\begin{align*}
R_{85}={}& -x_{5}x_{16} + x_{16} - x_{16}x_{97} + x_{99}
\end{align*}
% residual 86
\begin{align*}
R_{86}={}& -9x_{9} - 9x_{10} - 4x_{20} - 2x_{23} - 4x_{25} - 2x_{28} - 2x_{30} - 4x_{31} - 2x_{33} - 2x_{35} - 4x_{38} \\
&{} - 2x_{39} - 2x_{42} - 2x_{46} - 4x_{67} - 2x_{68} + x_{100}
\end{align*}
% residual 87
\begin{align*}
R_{87}={}& -x_{102}x_{105} + x_{105} + x_{107}
\end{align*}
% residual 88
\begin{align*}
R_{88}={}& -x_{4}x_{128} + x_{103}x_{210} - x_{106} - x_{107}x_{128}
\end{align*}
% residual 89
\begin{align*}
R_{89}={}& x_{108}x_{210} - x_{109}x_{131} - x_{110}
\end{align*}
% residual 90
\begin{align*}
R_{90}={}& x_{111}x_{210} - x_{112}x_{131} - x_{113}
\end{align*}
% residual 91
\begin{align*}
R_{91}={}& x_{114}x_{210} - x_{115}x_{131} - x_{116}
\end{align*}
% residual 92
\begin{align*}
R_{92}={}& x_{117}x_{210} - x_{118}x_{131} - x_{119}
\end{align*}
% residual 93
\begin{align*}
R_{93}={}& x_{120}x_{211} - x_{122}x_{131} - x_{123}
\end{align*}
% residual 94
\begin{align*}
R_{94}={}& x_{124}x_{211} - x_{125}x_{131} - x_{126}
\end{align*}
% residual 95
\begin{align*}
R_{95}={}& -x_{18}x_{128} + x_{131}
\end{align*}
% residual 96
\begin{align*}
R_{96}={}& x_{127} - x_{129}x_{131} - x_{130}
\end{align*}
% residual 97
\begin{align*}
R_{97}={}& -x_{131}x_{133} + x_{132} - x_{134}
\end{align*}
% residual 98
\begin{align*}
R_{98}={}& -x_{131}x_{136} + x_{135}x_{210} - x_{137}
\end{align*}
% residual 99
\begin{align*}
R_{99}={}& -x_{131}x_{139} + x_{138}x_{210} - x_{140}
\end{align*}
% residual 100
\begin{align*}
R_{100}={}& x_{141} - 2x_{141}x_{143} + x_{144}
\end{align*}
% residual 101
\begin{align*}
R_{101}={}& -x_{141}x_{144} - x_{144} + x_{145}
\end{align*}
% residual 102
\begin{align*}
R_{102}={}& 2x_{142} + x_{142}^{2} + x_{144}^{2} - x_{145}^{2}
\end{align*}
% residual 103
\begin{align*}
R_{103}={}& -x_{142}x_{147} - x_{147} + x_{148}
\end{align*}
% residual 104
\begin{align*}
R_{104}={}& -1 + x_{146}^{2} + 2x_{146}x_{148} + x_{147}^{2}
\end{align*}
% residual 105
\begin{align*}
R_{105}={}& x_{3}^{2} - 2x_{3}x_{142} - 2x_{142} + x_{149}
\end{align*}
% residual 106
\begin{align*}
R_{106}={}& -x_{3}x_{147} + x_{5} - x_{146} - x_{147} + x_{149}x_{150}
\end{align*}
% residual 107
\begin{align*}
R_{107}={}& 1 - 2x_{4} + x_{4}^{2} - 2x_{4}x_{142} + x_{151}
\end{align*}
% residual 108
\begin{align*}
R_{108}={}& -x_{4}x_{183} + x_{18} + x_{151}x_{152} - x_{182}
\end{align*}
% residual 109
\begin{align*}
R_{109}={}& -1 + 2x_{142}x_{194} - x_{153} + 2x_{194} - x_{194}^{2}
\end{align*}
% residual 110
\begin{align*}
R_{110}={}& -1 + x_{102} + x_{153}x_{154} - x_{182} - x_{183}x_{194}
\end{align*}
% residual 111
\begin{align*}
R_{111}={}& -x_{142}x_{156} - x_{156} + x_{157}
\end{align*}
% residual 112
\begin{align*}
R_{112}={}& -1 + x_{155}^{2} + 2x_{155}x_{157} + x_{156}^{2}
\end{align*}
% residual 113
\begin{align*}
R_{113}={}& 1 + x_{142} - x_{142}x_{159} - x_{159} + x_{160}
\end{align*}
% residual 114
\begin{align*}
R_{114}={}& x_{158}^{2} + 2x_{158}x_{160} - 2x_{159} + x_{159}^{2}
\end{align*}
% residual 115
\begin{align*}
R_{115}={}& -x_{142}x_{162} - x_{162} + x_{163}
\end{align*}
% residual 116
\begin{align*}
R_{116}={}& -1 + x_{161}^{2} + 2x_{161}x_{163} + x_{162}^{2}
\end{align*}
% residual 117
\begin{align*}
R_{117}={}& -x_{142}x_{165} - x_{165} + x_{166}
\end{align*}
% residual 118
\begin{align*}
R_{118}={}& -1 + x_{164}^{2} + 2x_{164}x_{166} + x_{165}^{2}
\end{align*}
% residual 119
\begin{align*}
R_{119}={}& -x_{142}x_{168} - x_{168} + x_{169}
\end{align*}
% residual 120
\begin{align*}
R_{120}={}& -1 + x_{167}^{2} + 2x_{167}x_{169} + x_{168}^{2}
\end{align*}
% residual 121
\begin{align*}
R_{121}={}& -x_{142}x_{171} - x_{171} + x_{172}
\end{align*}
% residual 122
\begin{align*}
R_{122}={}& -1 + x_{170}^{2} + 2x_{170}x_{172} + x_{171}^{2}
\end{align*}
% residual 123
\begin{align*}
R_{123}={}& -x_{142}x_{174} - x_{174} + x_{175}
\end{align*}
% residual 124
\begin{align*}
R_{124}={}& -1 + x_{173}^{2} + 2x_{173}x_{175} + x_{174}^{2}
\end{align*}
% residual 125
\begin{align*}
R_{125}={}& -x_{142}x_{177} - x_{177} + x_{178}
\end{align*}
% residual 126
\begin{align*}
R_{126}={}& -1 + x_{176}^{2} + 2x_{176}x_{178} + x_{177}^{2}
\end{align*}
% residual 127
\begin{align*}
R_{127}={}& -x_{142}x_{180} - x_{180} + x_{181}
\end{align*}
% residual 128
\begin{align*}
R_{128}={}& -1 + x_{179}^{2} + 2x_{179}x_{181} + x_{180}^{2}
\end{align*}
% residual 129
\begin{align*}
R_{129}={}& 1 - x_{142}x_{183} - x_{183} + x_{184}
\end{align*}
% residual 130
\begin{align*}
R_{130}={}& -1 + 2x_{182} + x_{182}^{2} + 2x_{182}x_{184} + x_{183}^{2}
\end{align*}
% residual 131
\begin{align*}
R_{131}={}& -x_{142}x_{186} - x_{186} + x_{187}
\end{align*}
% residual 132
\begin{align*}
R_{132}={}& -1 + x_{185}^{2} + 2x_{185}x_{187} + x_{186}^{2}
\end{align*}
% residual 133
\begin{align*}
R_{133}={}& -x_{142}x_{189} - x_{189} + x_{190}
\end{align*}
% residual 134
\begin{align*}
R_{134}={}& -1 + x_{188}^{2} + 2x_{188}x_{190} + x_{189}^{2}
\end{align*}
% residual 135
\begin{align*}
R_{135}={}& -x_{142}x_{192} - x_{192} + x_{193}
\end{align*}
% residual 136
\begin{align*}
R_{136}={}& -1 + x_{191}^{2} + 2x_{191}x_{193} + x_{192}^{2}
\end{align*}
% residual 137
\begin{align*}
R_{137}={}& -6 - 2x_{0} - 3x_{0}^{2} - 2x_{0}x_{3} - x_{1}^{2} - 2x_{3} - 2x_{3}^{2} - x_{4}^{2} - 2x_{5}^{2} + 2x_{9} - x_{9}^{2} - x_{10}^{2} \\
&{} - x_{11}^{2} - x_{12}^{2} - x_{13}^{2} - x_{14}^{2} - x_{15}^{2} - x_{16}^{2} - x_{17}^{2} - x_{18}^{2} - x_{100}^{2} + 2x_{102} - 2x_{102}^{2} \\
&{} - x_{103}^{2} - x_{104}^{2} - x_{108}^{2} - x_{111}^{2} - x_{114}^{2} - x_{117}^{2} - x_{120}^{2} - x_{121}^{2} - x_{124}^{2} - x_{127}^{2} \\
&{} - x_{128}^{2} - x_{132}^{2} - x_{135}^{2} - x_{138}^{2} + x_{141} - x_{194}^{2}
\end{align*}
% residual 138
\begin{align*}
R_{138}={}& x_{4} - 2x_{4}x_{9} - 2x_{4}x_{10} - 2x_{4}x_{11} - 2x_{4}x_{12} - 2x_{4}x_{13} - 2x_{4}x_{14} - 2x_{4}x_{15} \\
&{} - 2x_{4}x_{16} - 2x_{4}x_{17} + x_{194}
\end{align*}
% residual 139
\begin{align*}
R_{139}={}& -16 + x_{0} - x_{195}^{2}
\end{align*}
% residual 140
\begin{align*}
R_{140}={}& -x_{142}x_{197} - x_{197} + x_{206}
\end{align*}
% residual 141
\begin{align*}
R_{141}={}& -1 + x_{196}^{2} + x_{197}^{2} - x_{206}^{2}
\end{align*}
% residual 142
\begin{align*}
R_{142}={}& -x_{199} - x_{199}x_{200} + x_{207}
\end{align*}
% residual 143
\begin{align*}
R_{143}={}& -1 + x_{198}^{2} + x_{199}^{2} - x_{207}^{2}
\end{align*}
% residual 144
\begin{align*}
R_{144}={}& -x_{192}^{2} + x_{208}
\end{align*}
% residual 145
\begin{align*}
R_{145}={}& x_{197} - 2x_{202}x_{208} - x_{208}
\end{align*}
% residual 146
\begin{align*}
R_{146}={}& -4x_{197} - 8x_{197}x_{201} + x_{200}
\end{align*}
% residual 147
\begin{align*}
R_{147}={}& -x_{142} - x_{196}x_{203} + x_{200}
\end{align*}
% residual 148
\begin{align*}
R_{148}={}& -x_{191} - x_{193} - x_{196}x_{204} + x_{198}
\end{align*}
% residual 149
\begin{align*}
R_{149}={}& -x_{100} - 4x_{192}x_{205} + x_{199}
\end{align*}
% residual 150
\begin{align*}
R_{150}={}& 26 + 2x_{0}^{2} + 2x_{0}x_{3} - 2x_{0}x_{147} - 2x_{0}x_{156} + x_{1}^{2} - 2x_{1}x_{183} + x_{3}^{2} - 2x_{3}x_{156} + x_{5}^{2} \\
&{} - 2x_{5}x_{174} + x_{9}^{2} - 2x_{9}x_{159} + x_{10}^{2} - 2x_{10}x_{162} + x_{11}^{2} - 2x_{11}x_{165} + x_{12}^{2} \\
&{} - 2x_{12}x_{168} + x_{13}^{2} - 2x_{13}x_{171} + x_{14}^{2} - 2x_{14}x_{177} + x_{15}^{2} - 2x_{15}x_{180} + x_{16}^{2} \\
&{} - 2x_{16}x_{186} + x_{17}^{2} - 2x_{17}x_{189} + x_{100}^{2} - 2x_{100}x_{192} - x_{142}x_{209} - 2x_{146} - 2x_{148} \\
&{} + x_{148}^{2} - 2x_{155} - 2x_{157} + x_{157}^{2} - 2x_{158} + 2x_{159} - 2x_{160} + x_{160}^{2} - 2x_{161} - 2x_{163} \\
&{} + x_{163}^{2} - 2x_{164} - 2x_{166} + x_{166}^{2} - 2x_{167} - 2x_{169} + x_{169}^{2} - 2x_{170} - 2x_{172} + x_{172}^{2} \\
&{} - 2x_{173} - 2x_{175} + x_{175}^{2} - 2x_{176} - 2x_{178} + x_{178}^{2} - 2x_{179} - 2x_{181} + x_{181}^{2} - 2x_{182} \\
&{} + x_{184}^{2} - 2x_{185} - 2x_{187} + x_{187}^{2} - 2x_{188} - 2x_{190} + x_{190}^{2} - 2x_{191} - 2x_{193} + x_{193}^{2}
\end{align*}
% residual 151
\begin{align*}
R_{151}={}& -x_{104}x_{210} + x_{128}
\end{align*}
% residual 152
\begin{align*}
R_{152}={}& -x_{121}x_{211} + x_{128}
\end{align*}
% residual 153
\begin{align*}
R_{153}={}& -2 - x_{100}^{2} - x_{147}^{2} - x_{156}^{2} + 2x_{159} - x_{159}^{2} - x_{162}^{2} - x_{165}^{2} - x_{168}^{2} - x_{171}^{2} - x_{174}^{2} \\
&{} - x_{177}^{2} - x_{180}^{2} - x_{183}^{2} - x_{186}^{2} - x_{189}^{2} + x_{192}^{2} - x_{212}^{2} - x_{213}^{2} - x_{214} - x_{214}^{2}
\end{align*}
% residual 154
\begin{align*}
R_{154}={}& -1 - 4x_{106}^{2} - 4x_{110}^{2} - 4x_{113}^{2} - 4x_{116}^{2} - 4x_{119}^{2} - 4x_{123}^{2} - 4x_{126}^{2} + x_{128}^{2} - 4x_{130}^{2} \\
&{} - 4x_{134}^{2} - 4x_{137}^{2} - 4x_{140}^{2} - x_{215}^{2} - x_{216}^{2} - x_{217} - x_{217}^{2} - 4x_{219}^{2}
\end{align*}
% residual 155
\begin{align*}
R_{155}={}& -x_{4}x_{218} + x_{128}
\end{align*}
% residual 156
\begin{align*}
R_{156}={}& x_{2}x_{128} - x_{5}x_{218} - x_{219}
\end{align*}
% residual 157
\begin{align*}
R_{157}={}& -x_{2}x_{13} + x_{220}
\end{align*}
% residual 158
\begin{align*}
R_{158}={}& x_{102}x_{222} - 2x_{103}x_{156} + 2x_{108} - 2x_{108}x_{159} - 2x_{111}x_{162} - 2x_{114}x_{165} - 2x_{117}x_{168} \\
&{} - 2x_{120}x_{171} - 2x_{124}x_{177} - 2x_{127}x_{180} - 2x_{132}x_{186} - 2x_{135}x_{189} - 2x_{138}x_{192} \\
&{} - 2x_{142}x_{222} + 2x_{155}x_{156} + 2x_{156}x_{157} - 2x_{158} + 2x_{158}x_{159} + 2x_{159}x_{160} - 2x_{160} \\
&{} + 2x_{161}x_{162} + 2x_{162}x_{163} + 2x_{164}x_{165} + 2x_{165}x_{166} + 2x_{167}x_{168} + 2x_{168}x_{169} \\
&{} + 2x_{170}x_{171} + 2x_{171}x_{172} + 2x_{176}x_{177} + 2x_{177}x_{178} + 2x_{179}x_{180} + 2x_{180}x_{181} \\
&{} + 2x_{185}x_{186} + 2x_{186}x_{187} + 2x_{188}x_{189} + 2x_{189}x_{190} + 2x_{191}x_{192} + 2x_{192}x_{193} + x_{221} \\
&{} - 2x_{222}
\end{align*}
% residual 159
\begin{align*}
R_{159}={}& -1 - x_{102}x_{221} + x_{103}^{2} - 2x_{103}x_{155} + x_{108}^{2} - 2x_{108}x_{158} + x_{111}^{2} - 2x_{111}x_{161} \\
&{} + x_{114}^{2} - 2x_{114}x_{164} + x_{117}^{2} - 2x_{117}x_{167} + x_{120}^{2} - 2x_{120}x_{170} + x_{124}^{2} - 2x_{124}x_{176} \\
&{} + x_{127}^{2} - 2x_{127}x_{179} + x_{132}^{2} - 2x_{132}x_{185} + x_{135}^{2} - 2x_{135}x_{188} + x_{138}^{2} - 2x_{138}x_{191} \\
&{} + x_{155}^{2} - x_{156}^{2} + x_{158}^{2} + 2x_{159} - x_{159}^{2} + x_{161}^{2} - x_{162}^{2} + x_{164}^{2} - x_{165}^{2} + x_{167}^{2} \\
&{} - x_{168}^{2} + x_{170}^{2} - x_{171}^{2} + x_{176}^{2} - x_{177}^{2} + x_{179}^{2} - x_{180}^{2} + x_{185}^{2} - x_{186}^{2} + x_{188}^{2} \\
&{} - x_{189}^{2} + x_{191}^{2} - x_{192}^{2} + x_{222}
\end{align*}
% residual 160
\begin{align*}
R_{160}={}& x_{102}x_{224} - 2x_{104}x_{147} - 2x_{121}x_{174} - 2x_{128}x_{183} - 2x_{142}x_{224} + 2x_{146}x_{147} \\
&{} + 2x_{147}x_{148} + 2x_{173}x_{174} + 2x_{174}x_{175} + 2x_{182}x_{183} + 2x_{183} + 2x_{183}x_{184} + x_{223} \\
&{} - 3x_{224}
\end{align*}
% residual 161
\begin{align*}
R_{161}={}& -x_{102}x_{223} + x_{104}^{2} - 2x_{104}x_{146} + x_{121}^{2} - 2x_{121}x_{173} + x_{128}^{2} - 2x_{128}x_{182} + x_{146}^{2} \\
&{} - x_{147}^{2} + x_{173}^{2} - x_{174}^{2} + x_{182}^{2} - x_{183}^{2} + x_{223} + x_{224}
\end{align*}
\endgroup


\clearpage
\section{Artifact identities}
The repository snapshot is
\begin{center}\small\texttt{a0ef2df5bb1f73c35cd334bbc7f53553350ef06a}.\end{center}
The SHA-256 identities below fix the literal input and output.
\begin{description}[style=nextline]
\item[Pinned 245-variable parent]
{\small\texttt{1a42dd3a148c3084cd95a367869d251ed92e2f9852397215e4909e70b03d278e}}
\item[Final 225-variable residual specification]
{\small\texttt{fee84c408828aba3e7ceebe9a37f58509327063aa7fd078fd49c80adf3af4e56}}
\item[Collected 7,867-term polynomial]
{\small\texttt{3a9f582c1be2c4e4f927b383ffa66eb831a5a018c1f5d8441a16f58251ae14b0}}
\item[Final metrics]
{\small\texttt{b3f4606cf084a48a0652f3c3130e43adefecc8b76df1f4ae3a150033ee9cc054}}
\end{description}
The package-wide \path{SHA256SUMS} authenticates the remaining files,
including this report after PDF compilation. \path{provenance/SOURCES.json}
records the source origins and exact paths. Hashes identify artifacts;
their agreement does not replace the mathematical proof.

\end{document}
